\chapter{Elementos básicos en {\tt Mathematica}}

En este capítulo presentamos algunas de las funciones y comandos
básicos que se utilizan en el programa \verb"Wolfram Mathematica", que denominaremos simplemente
como \verb"Mathematica". Es de anotar que toda instrucción hecha en este programa se debe \emph{ejecutar} para obtener el producto señalado. En la versión \verb"13", bajo el sistema operativo \verb"Windows 10" esto se consigue con \verb"shift+enter" o \verb"enter" con el teclado numérico.

En general la sintaxis de las funciones inician con letra mayúscula
y en algunos casos llevan otras letras intermedias también en mayúscula. Los argumentos de cada función se insertan entre
corchetes \verb"[...]", los paréntesis \verb"(...)" se usan para agrupar y las llaves \verb"{...}" se utilizan para insertar \emph{listas}.

\section{Funciones básicas}
Las operaciones básicas: adición, sustracción, producto y cociente, entre dos cantidades $x$ e $y$ se denotan respectivamente como \verb"x+y", \verb"x-y", \verb"x*y" y \verb"x/y". 

\begin{tcolorbox}[breakable,enhanced,title=Advertencia]
Es importante aclarar que en general \verb"xy" no es igual a \verb"x*y".
\end{tcolorbox}
Algunas de las constante más utilizadas son: \verb"Pi", $\pi$; \verb"E", $e \approx 2.71828$;
\verb"I", la unidad imaginaria $i$, $i^{2}= -1;$ \verb"Infinity", $\infty$.

 Las funciones trigonométricas\footnote{Con argumento en radianes} se denotan como
{\small
\verb"Sin[x]," \verb"Cos[x]," \verb"Tan[x]," \verb"Cot[x]," 
\verb"Sec[x]," \verb"Csc[x]"
}
y sus respectivas funciones inversas
{\small
\verb"ArcSin[x]," \verb"ArcCos[x]," \verb"ArcTan[x]", \verb"ArcCot[x]," \verb"ArcSec[x]," \verb"ArcCsc[x]".
}

Para un punto $(x,y) \in \mathbb{R}^2$ se puede usar  \verb"ArcTan[x,y]", esta función posee como dominio a $[-\pi,\pi].$ 
Análogamente se escriben las funciones hiperbólicas y sus inversas \verb"Sinh[y]",  \verb"Cosh[y]", \verb"Tanh[y]", \verb"ArcSinh[y]", \verb"ArcCosh[y]," \verb"ArcTanh[y]", etc.

Existen tres tipos de igualdad, para asignación: \verb"x=2";  
prueba lógica: \verb"2+2==4", usada en ecuaciones y como una manera de probar dos expresiones son \emph{lo mismo}, 
\verb"x*x === x^2".
Otras funciones elementales son las siguientes.
\begin{tcolorbox}[breakable,enhanced,title=Algunas funciones básicas en Wolfram]
\begin{tabular}{rl}
\verb"x^y"& Potencia $y$ de $x$, $x^y$\\
\verb"Sqrt[x]"& Raíz cuadrada de $x$\\
\verb"Exp[x]"& Exponencial de $x$, $e^{x}$\\
\verb"Log[x]"& Logaritmo natural de $x$\\
\verb"Log[b,x]"& Logaritmo en base $b$ de $x$\\
\verb"n!"& Factorial de $n$\\
\verb"Abs[x]"& Valor absoluto de $x$\\
\verb"Max[x,y,...]", \verb"Min[x,y,...]"& Máximo o mínimo de $x, y,\cdots$\\
\verb"FactorInteger[n]"& Factores primos de $n$
\end{tabular}
\end{tcolorbox}

Para factorizar, simplificar, expandir, descomponer en fracciones parciales, operar fracciones o borrar una \verb"EXPRESIÓN"
se usan respectivamente, las siguientes funciones. \index{Together} \index{Simplify} \index{Apart}\index{Expand}
\begin{tcolorbox}[breakable, enhanced, title=Funciones básicas]
\begin{multicols}{2} % Activa las dos columnas
\begin{Verbatim}[formatcom=\letraverbatim]
Factor[EXPRESION]
Simplify[EXPRESION]
Expand[EXPRESION]
Apart[FRACCION]
Together[FRACCIONES]
Clear[VARIABLES]
\end{Verbatim}
\end{multicols}
\end{tcolorbox}

Por ejemplo, el usuario puede ejecutar las siguientes instrucciones:
\begin{tcolorbox}[breakable, enhanced]
\begin{multicols}{2}
\begin{Verbatim}[formatcom=\letraverbatim] 
Factor[x^16-y^16]
Apart[4*y^3/(y^4-16)]
Together[1/x+y/(y+1)]
Simplify[(x^9-1)/(x^3-1)]
Expand[(x^9-y^9)*(x^3-y^3)]
\end{Verbatim}
\end{multicols}
\end{tcolorbox}

Para extraer el numerador o el denominador de una expresión racional se pueden utilizar
las funciones \verb"Numerator" o \verb"Denominator". 
Utilizando las funciones \verb"ExpandDenominator" o \verb"ExpandNumerator"
se puede hacer la expansión en una expresión de los términos respectivos.
Las siguientes sentencias ilustran el uso de estos operadores.
\index{Numerator} \index{Denominator} \index{ExpandDenominator} \index{ExpandNumerator}
\begin{tcolorbox}[breakable,enhanced,title=Usos de \verb"ExpandDenominator" y \verb"ExpandNumerator"]
\begin{Verbatim}[formatcom=\letraverbatim]
Numerator[Together[x/y+y/x]]
Denominator[Together[x/y+y/x]]
ExpandDenominator[Together[y/(x+y)+x/(x-y)]]
ExpandNumerator[Together[y/(x+y)+x/(x-y)]]
\end{Verbatim}
\end{tcolorbox}

\index{TrigExpand} \index{TrigFactor} \index{TrigReduce}
Si una \verb"EXPRESIÓN" es una combinación de funciones trigonométricas, la instrucción
\verb"TrigExpand[EXPRESION]" la reescribe en términos de las funciones $\sin(\cdot)$ y $\cos(\cdot)$. 
La sentencia \verb"TrigFactor[EXPRESION]" factoriza usando funciones trigonométricas y
\verb"TrigReduce[EXPRESION]" convierte producto de funciones trigonométricas en suma. 
Por ejemplo:

\begin{tcolorbox}[breakable,enhanced]
\begin{multicols}{2}
\begin{Verbatim}[formatcom=\letraverbatim]
TrigExpand[Sin[5*x]-Cos[6*x]]
TrigReduce[Sin[2*x]*Cos[3*x]]
TrigFactor[Sin[5*x]-Cos[5*x]]
TrigFactor[Sin[x-y]-Cos[x+y]]
TrigExpand[Sin[x]^2*Cos[2*x]^3]
TrigReduce[Sin[3x]^3+Cos[2*x]^3]
\end{Verbatim}
\end{multicols}
\end{tcolorbox}

\subsection*{Tablas}
Para generar una lista con $n$ términos de la expresión $f(k)$ se usa la expresión \verb"Table[f[k],n]", además 
con \verb"Table[f[k],{k,a,b}]" se consigue que la lista satisfaga que $a \leq k \leq b.$ \index{Table}
Ilustremos esto con dos ejemplos sencillos.

\begin{tcolorbox}[breakable,enhanced]
\verb"Table[k^2,{k,2,20}"\\
\verb"Table[Prime[k],{k,2,20}"
\end{tcolorbox}

\subsection*{Vectores y funciones vectoriales}
El vector ${\bf u}= \langle a,b,c \rangle$ se define en \verb"Mathematica" como una lista \verb"u={a,b,c}".
El producto punto, el producto cruz y la norma de los vectores ${\bf u}$ y ${\bf v}$ se determina como sigue
\begin{tcolorbox}[breakable,enhanced,title=Vectores]
\verb"Dot[u,v]" \hskip1.75cm \verb"Cross[u,v]" \hskip1.75cm \verb"Norm[u]"
\end{tcolorbox}

\begin{tcolorbox}[breakable,enhanced,title=Funciones vectoriales]
Una función vectorial ${\bf r}(t)= \la x(t), y(t), z(t)\ra$ se define como un vector
\verb"r={x[t],y[t],z[t]}". 
\end{tcolorbox}
La instrucción  \verb"r={Cos[t],Sin[t],Cos[2*t]}" define una función vectorial que se denomina \verb"r".
Para extraer cada una de las tres componentes, 
se debe usar, según el caso, las siguientes instrucciones

%\verb"r[[1]]" \\ \verb"r[[2]]" \\ \verb"r[[3]]"
\begin{center}
\verb"r[[1]]" \hspace{3cm} \verb"r[[2]]" \hspace{3cm} \verb"r[[3]]"
\end{center}


Se pueden hacer algunos procedimientos asumiendo ciertas condiciones. Las funciones
\verb"Assuming" y \verb"Assumptions", permiten este tratamiento.
\index{Assumptions} \index{Assuming}
En los siguientes ejemplos se evidencia la diferencia de usar estas posibilidades.

\begin{tcolorbox}[breakable,enhanced,title=Usos de \verb"Assuming" y \verb"Assumptions"]
\begin{Verbatim}[formatcom=\letraverbatim]
Norm[{Cos[t],Sin[t],t}]
Norm[{Cosh[t],Sinh[t],t}]
Assuming[t>0,Simplify[Norm[{Cos[t],Sin[t],t}]]]
Assuming[t>0,Simplify[Norm[{Cosh[t],Sinh[t],t}]]]
Simplify[Norm[{Cosh[t],Sinh[t],t}],Assumptions->t>0]
\end{Verbatim}
\end{tcolorbox}

El símbolo $\%$ recupera el cálculo inmediatamente anterior, por ejemplo: 

\begin{tcolorbox}[breakable,enhanced]
\begin{multicols}{2}
\begin{Verbatim}[formatcom=\letraverbatim]
u={Exp[t],Exp[t],t}
v={-Exp[t],-Exp[t],t}
Norm[u+v]
Simplify[%,Assumptions->t>0]
\end{Verbatim}
\end{multicols}
\end{tcolorbox}

\subsection*{Sumas}
El comando \verb"Sum[Expresión,{k,a,b}]" obtiene la suma de una \verb"Expresión"
con variable de sumación $k$ la cual inicia en $a$ y termina en $b.$ 
Ejemplo de lo anterior son las siguientes sumas.
$$\sum_{k=1}^{10} k^2 \hskip1.5cm
\sum_{k=1}^n k^3 \hskip1.5cm
\sum_{k=1}^\infty \frac{1}{k^2} \hskip1.5cm
\sum_{r=0}^\infty \frac{2^r}{r!}$$
Para obtener su valor exacto se usan las siguientes instrucciones.

\begin{tcolorbox}[breakable, enhanced, title=Funciones \verb"Mathematica"]
\begin{multicols}{2} % Activa las dos columnas
\begin{Verbatim}[formatcom=\letraverbatim]
Sum[k^3,{k,1,n}]
Sum[k^2,{k,1,10}]
Sum[1/k^2,{k,1,Infinity}]
Sum[2^r/r!,{r,0,Infinity}]
\end{Verbatim}
\end{multicols}
\end{tcolorbox}
 
 
\subsection*{Límites}
Utilizando la sentencia \verb"Limit[F[x],x->a]" se evalúa en \verb"Mathematica" el $\lim_{x\to a} F(x)$.
También es posible calcular límites en varias variables como se ilustra enseguida.

{\small $$\lim_{x\to0}\frac{\sin(3x)}{\sin(2x)}\hskip1.5cm
\lim_{(x,y)\to (0,0)}\frac{x^2y}{x^2+y^2} \hskip1.5cm
\lim_{(x,y)\to(0,0)}\frac{\sin(x^2+y^2)}{x^2+y^2}$$}

Los límites anteriores se evalúan con las siguientes instrucciones.
\begin{tcolorbox}[breakable,enhanced]
\begin{Verbatim}[formatcom=\letraverbatim]
Limit[Sin[3*x]/Sin[2x],x->0]
Limit[x^2*y/(x^2+y^2),{x,y}->{0,0}]
Limit[Sin[x^2+y^2]/(x^2+y^2),{x,y}->{0,0}]
\end{Verbatim}
\end{tcolorbox}

\subsection*{Precaución}
Aunque es posible evaluar límites de manera iterada, es conveniente advertir que el uso 
incorrecto de esta sintaxis puede llevar a conclusiones arroneas. 
Por eso es importante usar la indicación adecuada. Ilustraremos esta situación.
\begin{itemize}%[leftmargin=0in]
\item Correcta \verb"Limit[n/(n+m),{n,m}->{Infinity,Infinity}]"
\item Incorrecta \verb"Limit[Limit[n/(n+m),m->Infinity],n->Infinity]"
\item Incorrecta \verb"Limit[Limit[n/(n+m),n->Infinity],m->Infinity]"
\end{itemize}

\subsection*{Operadores básicos}
%\begin{tcolorbox}[breakable,enhanced,title=Operadores Básicos en \verb"Mathematica"]
\index{Solve} \index{NSolve}
\begin{itemize}
\item Para resolver la ecuación $f(x)=0,$ en donde $x$ es la incognita, se usa la
instrucción: \verb"Solve[f[x]==0,x]". La solución numérica de ecuaciones se resuelve utilizando \verb"NSolve".
{\bf No olvidar $"=="$}.
\item La derivada de $f(x)$ con respecto a $x$, se obtiene con \verb"D[f[x],x]", 
la segunda derivada de $f(x)$ se consigue con \verb"D[f[x],x,x]".
La $n-$ésima derivada de $f(x)$ con respecto a $x$, $\frac{\partial^n f}{\partial x^n}$, se obtiene con la instrucción
\verb"D[f[x],{x,n}]". 
\item La derivada parcial de $f(x,y,z)$ con respecto a $x$, se consigue ejecutando \verb"D[f[x,y,z],x]", mientras que 
la diferencial total de $x^y,$ (suponiendo que $y$ depende de $x$) se obtiene mediante la expresión 
\verb"Dt[x^y,x]". \index{Diferencial total}
\item La integral indefinida de la función $f(t)$ con respecto a $t$, 
se evalúa mediante la instrucción: \hskip0.1cm \verb"Integrate[f(t),t]".
Con la instrucción \\\verb"Integrate[f(t),{t,a,b}]" se evalúa la integral definida $\ds \int_a^bf(t)dt$.
  \index{Integrate}
\end{itemize} 

%\end{tcolorbox} 
Ejemplo de lo anterior son las siguientes instrucciones:
\begin{tcolorbox}[breakable,enhanced]
\begin{multicols}{2}
\begin{Verbatim}[formatcom=\letraverbatim]
Solve[x^4==1,x]
Solve[x^4y^2==1,y,Reals]
D[ArcSin[x],x]
D[Sqrt[1/(1+r)], r]
D[Sqrt[x^3+x-2], x,x]
Integrate[1/(1+r),{r,0,1}]
\end{Verbatim}
%\end{itemize}
\end{multicols}
\end{tcolorbox}

\subsection*{Animación de objetos}
Con los siguientes comandos se puede dar \emph{animación} a un \verb"TERMINO", la variable es $k$, con $n \leq k \leq m $
e incrementos de $r$ unidades.
\index{Manipulate}\index{Animate}
\begin{tcolorbox}[breakable,enhanced,title=Las funciones {\tt Manipulate} y {\tt Animate} ]

\begin{Verbatim}
Manipulate[TERMINO,{k,n,m,r}]
Animate[TERMINO,{k,n,m,r}]
\end{Verbatim}

\end{tcolorbox}
La $r$ se puede omitir. % en el caso de la función \verb"Animate". 
El lector puede ejecutar las siguientes instrucciones
para evidenciar un poco mejor, el uso de las funciones mencionadas
\begin{tcolorbox}[breakable,enhanced,title=Ejemplos con {\tt Manipulate} y {\tt Animate}]
\begin{Verbatim}
Manipulate[Expand[(a+b)^k],{k,1,20,1}]
Animate[Expand[(a+b)^k],{k,1,20,1}]
Manipulate[Integrate[x^k,x],{k,-5,5,1/2}]
Manipulate[Integrate[(Sin[x])^k,x],{k,-2,10,1}]
\end{Verbatim}
\end{tcolorbox}

\subsection*{Transformación a coordenadas polares}
La función \verb"ToPolarCoordinates" transforma el punto $(a,b)$ escrito en coordenadas cartesianas 
al punto $(r, t)$ escrito en coordenadas polares. 
Con la función \verb"FromPolarCoordinates" se realiza la transformación inversa.
\index{ToPolarCoordinates}  \index{FromPolarCoordinates}

\begin{tcolorbox}[breakable,enhanced]
\verb"ToPolarCoordinates[{a,b}]"\\
\verb"FromPolarCoordinates[{r,t}]"
\end{tcolorbox}
 
El argumento $t$ satisface que $-\pi \leq t \leq \pi$. Por ejemplo:
\begin{tcolorbox}[breakable,enhanced]
\verb"ToPolarCoordinates[{-1,-1}]"\\
\verb"FromPolarCoordinates[{2,-3*Pi/4}]"
\end{tcolorbox}

\subsection*{Definición de funciones}
Es posible definir una función y luego usarla para cálculos posteriores,
esta puede depender de una o varias variables, la sintaxis básica es:

\begin{tcolorbox}[breakable,enhanced,title=Definición de funciones]
\verb"f[x_] := Expresión" para funciones de una variable $x.$ \\
\verb"f[x_,y_] := Expresión" para funciones de dos variables $x,y.$
\end{tcolorbox}
% Al no definir la función en un conjunto el valor por defecto es $0.$ 
Con la siguiente sentencia se define la función vectorial ${\bf r}(t,u)$, \\
\verb"r[t_,u_]:={(1-u)*Cos[2t]+u*Sin[3t],u*Cos[3t],t(1-u)}" \\
para evalúarla en $t=\pi/3$ y $u=1/2$, se utiliza \verb"r[Pi/3,1/2]".

Se pueden definir funciones a trozos, con $k$ condiciones y $x \in [a,b]$, 
\begin{tcolorbox}[breakable,enhanced,title=Definición de funciones a trozos]
\verb"Piecewise[{{Expresión 1,condición 1},{Expresión 2,"
\verb"condición 2},...,{Expresión k,condición k}},{x,a,b}]"
\end{tcolorbox}  \index{Piecewise}

\enlargethispage{\baselineskip}

Ahora se muestran algunos ejemplos:
\begin{itemize}[leftmargin=0cm,itemindent=.5cm,labelwidth=\itemindent,labelsep=0cm,align=left]
\item Para definir la función $f(x,y)=x^2-3y^4$ y evaluarla en el punto $(2,3)$ se ejecutan 
las instrucciones \verb"f[x_,y_]:=x^2-3*y^4", \verb"f[2,3]".
\item La función 
$$ f(x,y) = \left \{ \begin{array}{cl} \frac{xy}{x^2-y^2} & \text{ si } x>y \\[0.2cm] \sin(x+y) & \text{en caso contrario} \end{array} \right.$$
se define ejecutando la siguiente instrucción:\\
\verb"f[x_,y_]:=Piecewise[{{x*y/(x^2-y^2),x>y}},Sin[x+y]]"
\item Con la instrucción \verb"f[x_]:=Piecewise[{{4-x^2,x<=1}," \\
\verb"{2-x^3,1<x<2},{2*Sin[x],x>=2}}]", se define la función 
$$ f(x,y) = \left \{
 \begin{array}{cl} 4-x^2 & \text{ si } x\leq 1 \\[0.2cm]
2-x^3 & \text{ si } x\in(1,2)\\[0.2cm]
2 \sin(x) & \text{ si } x\geq 2
\end{array} \right.$$
\end{itemize}

\section{Gráficos}
Presentamos las principales posibilidades que ofrece \verb"Mathematica" 
para obtener la representación gráfica de objetos matemáticos.

\subsection*{Gráficas de funciones en el plano}
Se ofrece una manera de obtener la representación gráfica en el plano de objetos o funciones 
escritas en forma explícita, implícita o paramétrica.

\begin{itemize}[leftmargin=0cm,itemindent=.5cm,labelwidth=\itemindent,labelsep=0cm,align=left]
\item \verb"Plot[f[x],{x,a,b}]" dibuja la función $y=f(x)$ con $x \in[a,b].$ \index{Plot}
\item \verb"PolarPlot[r[t],{t,a,b}]" grafica la función $r(t)$ escrita en coordenadas polares, con $t \in[a,b].$
\item \verb"ContourPlot[f[x,y]==0,{x,a,b},{y,c,d}]",\index{ContourPlot}
dibuja la ecuación de dos variables $f(x,y)=0$, con $x \in[a,b]$ y $y \in[c,d],$
se debe incluir el símbolo de igualdad y no olvidar $"=="$.
\item \verb"ParametricPlot[{x[t],y[t]},{t,a,b}]", dibuja la curva parametri- \\ zada $\ve{r}(t)=\left\langle x(t),y(t)\right\rangle ,\,\, t\in \lbrack a,b].$
\index{ParametricPlot}
\end{itemize}
% \end{tcolorbox}
Si se desean incluir varias expresiones en las funciones anteriores,
se puede hacer una lista con estas y encerrarlas en llaves (\verb"{...}").
Las siguientes instrucciones son ejemplos de las anteriores funciones:

\begin{tcolorbox}[breakable,enhanced,title=Algunos gráficos en el plano]
\begin{Verbatim}
Plot[Exp[-10*t]*(2+20t)+1,{t,0,1}]
PolarPlot[1-Cos[t],{t,0,2*Pi}] 
PolarPlot[3/(6+Cos[10*u]),{u,0,2*Pi}]
ContourPlot[{y^2+(1-x^2)^2==4,y^2+x^2==4},{x,-2,2},
{y,-2,2}] 
\end{Verbatim}
\begin{Verbatim}
ParametricPlot[{Cos[t],Sin[2*t]},{t,0,2Pi}]
Manipulate[PolarPlot[t*Sin[k*t],{t,0,2Pi}],
{k,1,5,0.5}]
ParametricPlot[{{2*Cos[t],2*Sin[t]},{Cos[t],Sin[t]},
{Cos[t]*(1-Cos[t]),Sin[t]*(1-Cos[t])}},{t,0,2*Pi}]
\end{Verbatim}
\end{tcolorbox}

\subsection*{Gráficos de funciones en el espacio}
De manera análoga a las funciones descritas anteriormente para
conseguir gráficos en el plano, se puede agregar la extensión \verb"3D" y unas opciones adicionales
con el fin de obtener gráficos en el espacio.
Presentamos unas de las funciones más utilizadas en  \verb"Mathematica".
\\

\begin{itemize}[leftmargin=*]
\item \verb"Plot3D[f[x,y],{x,a,b},{y,c,d}]", permite obtener el gráfico de la función 
$z=f(x,y)$, para $(x,y) \in [a,b] \times [c,d]$. 
\\


\item \texttt{Plot3D[\allowbreak\{f[x,y],\allowbreak g[x,y],\allowbreak ..., h[x,y]\},\allowbreak \{x,a,b\},\allowbreak \{y,c,d\}]} \\
permite obtener el gráfico de varias funciones.
\\
%Mientras que con \verb"Plot3D[{f[x,y],g[x,y],...,h[x,y]},{x,a,b},{y,c,d}]", 
%se obtiene el gráfico de varias funciones. \index{Plot3D}
\item \Verb!ParametricPlot3D[{x[t],y[t],z[t]},{t,a,b}]!,
dibuja la curva parametrizada $\ve{r}(t)=\langle x(t),y(t),z(t) \rangle $, 
con $t\in[a,b].$  \index{ParametricPlot3D}
\\

\item \verb"ParametricPlot3D[{x[u,v],y[u,v],z[u,v]},{u,a,b},{v,c,d}]",
se utiliza para graficar superficies parametrizadas por medio de una función
$\ve{r}(u,v)=\langle x(u,v),y(u,v),z(u,v) \rangle $, con
$u\in[a,b],$ $v\in[c,d].$
\\
\item \verb"ContourPlot3D[f[x,y,z]==0,{x,a,b},{y,c,d},{z,p,q}]", se usa
para graficar conjuntos de nivel (contornos).
Este comando presenta el gráfico de una ecuación de tres variables, con $x\in[a,b],$
$y\in[c,d]$ y $z\in[p,q].$  \index{ContourPlot3D}
\\
\end{itemize}

Los gráficos se pueden personalizar en sus atributos con el uso de algunas opciones, 
presentamos las más utilizadas:

\vspace{-2pt}

\begin{tcolorbox}[breakable,enhanced,title=Opciones de los gráficos]
\begin{tabular}{rl}
\verb"Axes" & Muestra los ejes, \verb"True/False"\\
\verb"AxesLabel->{u,v,w}" & Etiqueta los ejes como $u,$ $v$ y $w$ \\
\verb"AxesOrigin->{a,b,c}" & Muestra los ejes con centro en $(a,b,c)$ \\
\verb"Blend[{color1,...}]"& Permite combinar varios colores \\
\verb"Boxed" & Muestra gráfico en caja, \verb"True/False"\\
\verb"BoxRatios->{a,b,c}" & Proporciones de la caja en un gráfico\\
\verb"ColorFunction" & Determina el color de la función\\
\verb"ContourStyle" & Opciones en \verb"ContourPlot3D" \\
\verb"Dashed" & Presenta líneas discontinuas\\
\verb"Mesh->{n,m}" & Crea la malla, con $n$ y $m$ líneas \\
\verb"Mesh" & Activa la malla, \verb"True/False" \\ 
\verb"MeshStyle" & Estilo de la malla (color)\\
\verb"Opacity[a]" &  Opacidad, se usa en plotStyle, $0 \leq a \leq 1$\\
\verb"PlotStyle" &  Opciones del gráfico, color, líneas, etc.\\
\verb"PlotRange" &  Rango de valores en el gráfico\\
\verb"PlotPoints" & Número de puntos sobre el dibujo\\
\verb"RGBColor[a,b,c]" & Crea un color con proporciones $a$, $b$, $c$\\
\verb"Thick" & Crea una línea más gruesa \\
\verb"Thickness" & Grosor de la línea \\
\verb"Ticks" &  Indica las marcas deseadas en la caja\\
\verb"Tube[a]" & Configura un tubo de radio $a$. \\
\end{tabular}
\end{tcolorbox}

El lector puede además de ejecutar las instrucciones, modificarlas.
{\small
\begin{itemize}
\item \verb"Plot3D[{5-y^2,x^2+y^2},{x,-4,4},{y,-4,4},PlotPoints->60,"\\
\verb"PlotStyle->Opacity[0.6]]"
\item \verb"ContourPlot3D[{z==6-(x-1)^2-y^2,4*x^2+y^2==36,"\\
\verb"x^2+4*y^2==4},{x,-6,10},{y,-7,7},{z,-35,6},"\\
\verb"ContourStyle->{Red,Yellow,Green}]"
\item \verb"ParametricPlot3D[{Cos[u]*Cosh[v],Sin[u]*Cosh[v],Sinh[v]},"\\
\verb"{u,0,2*Pi},{v,-2,2},PlotStyle->Orange,MeshStyle->Gray]"
\item \verb"ContourPlot3D[z==3*x-x^3-2*y^2+y^4,{x,-3,3},{y,-2,3},"\\
\verb"{z,-3,3},MeshStyle->Gray]"
\item \begin{Verbatim} 
ContourPlot3D[z==3*x-x^3-2*y^2+y^4,{x,-3,3},{y,-2,3},
{z,-3,3},ColorFunction->Function[{x,y,z},
ColorData["Rainbow"][z]],MeshStyle->Gray]
\end{Verbatim}
\end{itemize}
}

El siguiente gráfico se consigue utilizando las instrucciones señaladas.


\begin{figure}[H]
\centering
\caption{Integración en el espacio}
\includegraphics[width=0.4\textwidth]{Img/math_01a.pdf} \hspace{5mm}
\includegraphics[width=0.4\textwidth]{Img/math_01b.pdf}
\fuentepropia
\end{figure}

\begin{tcolorbox}[breakable,enhanced, title=Instrucción en \verb"Mathematica"]
\begin{Verbatim}[fontsize=\small]
ContourPlot3D[z==(y^2-x)*Exp[1-x^2-y^2],{x,-2,2},
{y,-2,2},{z,-1,2},Mesh->{6,8},Boxed->False,
MeshStyle->RGBColor[0.5,0.5,0.5],AxesOrigin->{0,0,0},
ContourStyle->Orange,Ticks->{{-1,1,2},{-1,1,2},
{-1,1}},AxesStyle->Black,AxesLabel->{u,v,w}]
\end{Verbatim}
\end{tcolorbox}


La composición de varias funciones puede hacer complicado
el tratamiento, la función \verb"Show" permite ver, en un solo gráfico, varios objetos obtenidos con 
diferentes funciones, por ejemplo: \index{Show}
\begin{tcolorbox}[breakable,enhanced,title=Uso de la función \verb"Show"]
\begin{Verbatim} 
Show[ParametricPlot3D[{Cos[t],Sin[t],1},{t,0,2Pi},
PlotStyle->Red],Graphics3D[{Blue,Arrowheads[0.05],
Arrow[Tube[{{1,0,1},{1,1,1}},0.02]]}],
AxesOrigin->{0,0,0},
PlotRange->{{-2,2},{-2,2},{-1,2}}]
\end{Verbatim}
\end{tcolorbox}
Otra posibilidad es ejecutar las instrucciones de manera separada y luego mostrarlos en un solo 
gráfico con la función \verb"Show". En el siguiente ejemplo, las primeras dos instrucciones se ejecutan, 
se nombran como \verb"A" y \verb"B", y luego con la función \verb"Show" se muestran ambas.%\\
\begin{tcolorbox}[breakable,enhanced,title=Ejecución de gráficos por separado]
\begin{Verbatim} 
A=ContourPlot3D[{x==4-y-4z+2z^2,(1-z)^2+(2-y)^2==1}
{x,-3,3},{y,-1,4},{z,-1/2,5/2},Mesh->False];
B=ParametricPlot3D[{Sin[t]+2(Cos[t])^2,2-Sin[t],
1-Cos[t]},{t,0,2Pi},
PlotStyle->{Red,Thickness->0.02}];
Show[A,B]
\end{Verbatim}
\end{tcolorbox}
Otras posibilidades muy útiles para obtener gráficos son las siguientes.

\subsection*{Gráfico de puntos, segmentos y vectores}

\begin{tcolorbox}[breakable,enhanced,title=Gráfico de puntos]
La instrucción \verb"Graphics[Point[{a,b}]]" dibuja el punto de
coordenadas $\left(a,b\right).$ Con \verb"Graphics3D[Point[{a,b,c}]]" se dibuja el 
punto de coordenadas $\left(a,b,c\right) .$
\end{tcolorbox}
El tamaño de un punto se modifica con el comando \verb"PointSize". \index{PointSize}
Para dibujar de color rojo los puntos $\left( 1,1,1\right) $ y $\left( 2,0,1\right)$
se ejecuta la instrucción \verb"Graphics3D[{Red,PointSize[0.1],Point[{{1,1,1},{2,0,1}}]}]".

El segmento que une el punto $(a,b)$ con el punto $(p,q)$ se representa utilizando la siguiente instrucción:
\begin{tcolorbox}[breakable,enhanced,title=Gráfico de segmentos]
\begin{Verbatim}
Graphics[Line[{{a,b},{p,q}}]]
\end{Verbatim}
\end{tcolorbox} \index{Line}

Con la instrucción \texttt{Graphics3D[\allowbreak Line[\allowbreak\{\{a,b,c\},\allowbreak\{p,q,r\}\}]]}, se obtiene el gráfico de un segmento tridimensional. Por ejemplo, con la instrucción \texttt{Graphics3D[\allowbreak\{\allowbreak Dashed,\allowbreak Blue,\allowbreak Line[\allowbreak\{\{1,-1,1\},\allowbreak\{2,3,-1\}\}]\}]}, se obtiene el segmento de color azul en línea discontinua que une los puntos $(1,-1,1)$ y $(2,3,-1).$

También es posible dibujar el vector que va del punto $(a,b,c)$ al punto $(p,q,r).$
\begin{tcolorbox}[breakable,enhanced,title=Gráfico de vectores en el espacio]
\begin{Verbatim}
Graphics3D[Arrow[{{a,b,c},{p,q,r}}]]
\end{Verbatim}
\end{tcolorbox}
Si se desea agregar determinado grosor $u$ al vector, la instrucción que se debe ejecutar es 
\verb"Graphics3D[Arrow[{{a,b,c},{p,q,r}},u]]". \index{Graphics3D}
El vector que une el punto $(1,1,-1)$ con el punto $(2,2,0),$ y algunos atributos personalizados, 
se dibuja en  utilizando las siguientes instrucciones:
\begin{tcolorbox}[breakable,enhanced]
\begin{Verbatim}[formatcom=\letraverbatim]
Graphics3D[{Black,Arrowheads[0.1],
Arrow[Tube[{{1,1,-1},{2,2,0}},0.01]]}]
\end{Verbatim}
\end{tcolorbox}

\begin{figure}[H]
	\begin{center}
		\caption{El vector $\langle 1,1,1\rangle$ se dibuja en el punto $(1,1,-1)$.}
		\includegraphics[height=4.5cm]{Img/T_67a.pdf} \hskip0.2cm
		\includegraphics[height=4.5cm]{Img/T_67b.pdf}
\fuentepropia
	\end{center}  
\end{figure}

\subsection*{Regiones en el plano y el espacio}
La función \verb"RegionPlot" muestra el gráfico de una región dimensional, la sintaxis puede incluir desigualdades 
y los operadores lógicos: \emph{y} (\verb"&&"), \emph{o} (\verb"||"). En los siguientes ejemplos se incorpora el uso 
de las opciones \verb"BoundaryStyle" y \verb"Thickness" que controla el grosor y color de la frontera.
\index{Operadores lógicos} \index{BoundaryStyle} \index{Thickness} 


\begin{figure}[H]
\centering
\caption{\small Uso de la función {\tt RegionPlot}}
\includegraphics[width=0.4\textwidth]{Img/T_49a.pdf} \hspace{5mm}
\includegraphics[width=0.4\textwidth]{Img/T_49b.pdf}

\includegraphics[width=0.4\textwidth]{Img/T_49c.pdf} \hspace{5mm}
\includegraphics[width=0.4\textwidth]{Img/T_49d.pdf}
\fuentepropia
\end{figure}

La versión tridimensional es \verb"RegionPlot3D".



\begin{tcolorbox}[breakable,enhanced]
\begin{Verbatim}[formatcom=\letraverbatim] 
RegionPlot[1<=x^2+2*y^2<=9&&y>=-x^2,{x,-3,3},
{y,-2,5/2},PlotPoints->50,BoundaryStyle->{Black,
Thickness->0.01},PlotStyle->Gray,Frame->False]
RegionPlot[(x+1)^2+2*y^2<3||(x-1)^2+y^2<3,{x,-3,3},
{y,-2,2},BoundaryStyle->{Black,Thickness->0.01},
PlotStyle->Gray,Frame->False]
RegionPlot[{(x+1)^2+2*y^2<3,(x-1)^2+y^2<3},{x,-3,3},
{y,2,2},PlotStyle->Gray,BoundaryStyle->{Black,
Thickness->0.01},Frame->False]
RegionPlot[x^2+y^2<9&&Abs[x]+Abs[y]>2,{x,-3,3},
{y,-3,3},Frame->False,PlotStyle->Gray,
BoundaryStyle->{Black,Thickness->0.01}]
\end{Verbatim}
\end{tcolorbox}



Para representar un sólido se utiliza la función \verb"RegionPlot3D". \index{RegionPlot3D}
Este comando dibuja la región a partir de ciertas \verb"DESIGUALADES", con $x\in[a,b]$, $y\in[c,d]$, $z\in[u,v].$ 
Las inecuaciones se concatenan con los símbolos \verb"&&".
La sintaxis básica es:



\begin{tcolorbox}[breakable,enhanced,title=Gráfico de una región sólida]
\begin{Verbatim}[formatcom=\letraverbatim]
RegionPlot3D[DESIGUALDADES,{x,a,b},{y,c,d},{z,u,v}]
\end{Verbatim}
\end{tcolorbox}


Presentamos ahora, algunas situaciones que ilustran el uso de esta función.

(1) El sólido acotado por las figuras de $z=0,$ $z=1-y,$ $x^2=y,$ $y=1.$


\begin{figure}[H]
\centering
\caption{Cono truncado oblicuo determinado por ${\bf r}(u,t)$}
\includegraphics[width=0.4\textwidth]{Img/math_06c.pdf} \hspace{5mm}
\includegraphics[width=0.4\textwidth]{Img/math_06d.pdf}
\fuentepropia
\end{figure}

Esta región se obtiene con las siguientes instrucciones:

\begin{tcolorbox}
\begin{Verbatim}[formatcom=\letraverbatim]
RegionPlot3D[0<z<1-y&&x^2<y<1,{x,-1,1},{y,0,1},
{z,0,1},PlotPoints->90,PlotStyle->Orange,
Mesh->False,Ticks->{{-1,0,1},{-1,1},{0,1}}]
\end{Verbatim}
\end{tcolorbox}

%\clearpage
(2) El sólido acotado por las figuras de $(x-1)^2+y^2= 1,$ $(x-3/2)^2+y^2=4,$ $z=0,$ $2z+y=4.$



\begin{figure}[H]
\centering
\caption{\small Sólido acotado por $x^2+2x+y^2=0,$ $x^2+y^2-3x=\frac{7}{4},$ $2z+y=4.$}
\includegraphics[width=0.4\textwidth]{Img/math_06.pdf} \hspace{5mm}
\includegraphics[width=0.4\textwidth]{Img/math_06a.pdf}

\includegraphics[width=0.4\textwidth]{Img/math_06b.pdf}
\fuentepropia
\end{figure}

Para obtenerlo se ejecutó la siguiente instrucción:
%
\begin{tcolorbox}
\begin{Verbatim}[formatcom=\letraverbatim] 
RegionPlot3D[z>0&&z<4-y/2&&(x-1)^2+y^2>1&&(x-1.5)^2
+y^2<4,{x,-1,4},{y,-2,2},{z,0,5},Mesh->False,
PlotPoints->90]
\end{Verbatim}
\end{tcolorbox}
%
(3) Las expresiones $(x^2+y^2)^2=z+1, \,\,\, z= 4(1-x^2-y^2), $ 
acotan el sólido que se muestra enseguida. 

\begin{figure}[H]
	\begin{center}
		\caption{Superficie determinada por $(x^2+y^2)^2=z+1$ y $z= 4(1-x^2-y^2)$.}
		\includegraphics[height=4.5cm]{Img/math_06e.pdf} \hspace{0.3cm}
		\includegraphics[height=4.5cm]{Img/math_06f.pdf}
\fuentepropia

\end{center}
\end{figure}

Las siguientes instrucciones permiten obtenerlo:
\begin{tcolorbox}
\begin{Verbatim}[formatcom=\letraverbatim]
RegionPlot3D[(x^2+y^2)^2-1<z<4*(1-x^2-y^2)^2&&
-Sqrt[1-x^2]<y<Sqrt[1-x^2],{x,-1,1},{y,-1,1},
{z,-1,4},PlotPoints->90,Mesh->False]
\end{Verbatim}
\end{tcolorbox}

(4) La región del espacio interior a la esfera $x^2+y^2+z^2=4$ y al hiperboloide de dos mantos $x^2+y^2=2+z^2$, exterior al cilindro $x^2+y^2=\frac{2}{5}$ se muestra en el siguiente gráfico.


\begin{figure}[H]
\centering
\caption{\small Sólido acotado por $x^2+y^2+z^2=4,$ $x^2+y^2=2+z^2,$ $x^2+y^2=\frac{2}{5}.$}
\includegraphics[width=0.4\textwidth]{Img/math_06n.pdf} \hspace{5mm}
\includegraphics[width=0.4\textwidth]{Img/math_06p.pdf}
\fuentepropia
\end{figure}

Este se consigue con las siguientes instrucciones:

\begin{tcolorbox}
\begin{Verbatim}[formatcom=\letraverbatim]
RegionPlot3D[x^2+y^2+z^2<4&&x^2+y^2<2+z^2&&
x^2+y^2>0.4,{x,-2,2},{y,-2,2},{z,-2,2},PlotPoints->90]
\end{Verbatim}
\end{tcolorbox}


(5) En coordenadas cilíndricas, las desigualdades 
$1 \leq r \leq 2,\,\, r \leq z \leq 2, $
determinan la región acotada por un cilindro, un semicono y un plano.


\begin{figure}[H]
\centering
\caption{\small Sólido acotado por $1 \leq r \leq 2,$ $r \leq z \leq 2.$}
\includegraphics[width=0.4\textwidth]{Img/math_06g.pdf} \hspace{5mm}
\includegraphics[width=0.4\textwidth]{Img/math_06h.pdf}
\end{figure}

Para obtener estas figuras se definen $r$ y $\theta$, 
luego se escriben las desigualdades y se agregan otras opciones:
\begin{tcolorbox} 
\begin{Verbatim}[formatcom=\letraverbatim]
{r,\[Theta]}={Sqrt[x^2+y^2],ArcTan[y/Abs[x]]};
RegionPlot3D[1<=r<=2&&r<=z<=2,{x,-2,2},{y,-2,2},
{z,1,2},PlotPoints->90,Mesh->True,AxesLabel->
Automatic,PlotRange->{{-2,2},{-2,2},{1,2.2}}]
\end{Verbatim}
\end{tcolorbox}


\subsection*{Gráficos en coordenadas esféricas}
En coordenadas esféricas se puede utilizar la 
función \verb"SphericalPlot3D", con una sintaxis genérica para graficar una función de la forma 
$\rho = f(\theta,\phi)$, con  $\theta \in [a,b]$ y $\phi \in [c,d];$
esto es: \index{SphericalPlot3D}

\begin{tcolorbox}[breakable,enhanced,title=Gráficos en coordenadas esféricas]
\begin{lstlisting}[language=Mathematica, basicstyle=\ttfamily]
SphericalPlot3D$[f[\[Theta],\[Phi]], 
{\[Theta],a,b}, {\[Phi],c,d}]
\end{lstlisting}
\end{tcolorbox}

Así, se ejecutan las siguientes instrucciones:

\begin{tcolorbox}[breakable,enhanced]
\begin{Verbatim}[formatcom=\letraverbatim]
SphericalPlot3D[1,{\[Theta],Pi/6,Pi},{\[Phi],
Pi/6,5*Pi/3}]
SphericalPlot3D[Cos[\[Phi]],{\[Theta],0,Pi},{\[Phi],
0,Pi}]
SphericalPlot3D[\[Theta],{\[Theta],0,Pi},{\[Phi],
0,2*Pi}]
\end{Verbatim}
\end{tcolorbox}



Es así como se obtienen, respectivamente, los resultados de la figura 21:


\begin{figure}[H]
\centering
\caption{\small Regiones en coordenadas esféricas.}
\label{regiones_esfe}
\includegraphics[width=0.4\textwidth]{Img/math_06aa.pdf} \hspace{5mm}
\includegraphics[width=0.4\textwidth]{Img/math_06ab.pdf}

\includegraphics[width=0.4\textwidth]{Img/math_06ac.pdf}
\end{figure}

De manera alternativa, se pueden usar desigualdades definiendo previamente las variables $\rho$, $\theta$ y $\phi.$ 
Por ejemplo, la región limitada por las siguientes expresiones
$1 \leq \rho \leq 2,\,\, 0\leq \theta\leq \frac{3\pi}{4}, \,\,
\frac{\pi}{6} \leq \phi \leq \frac{3\pi}{4}, $
se muestra a continuación. Para conseguirlo, se pueden definir las variables $\rho,\,\, \theta,\,\, \phi$
y luego establecer la región que se ilustra a continuación:


\begin{figure}[H]
\centering
\caption{\small Región determinada por $1\leq \rho \leq 2,$ $0\leq \theta\leq \frac{3\pi}{4}$ y $\frac{\pi}{6} \leq \phi \leq \frac{3\pi}{4}$}
\includegraphics[width=0.4\textwidth]{Img/math_06m.pdf} \hspace{5mm}
\includegraphics[width=0.4\textwidth]{Img/math_06l.pdf}

\includegraphics[width=0.4\textwidth]{Img/math_06k.pdf}
\fuentepropia
\end{figure}

Para conseguirlo, se pueden definir las variables $\rho$, $\theta$, $\phi$ y luego establecer la región que se ilustra a continuación:

\begin{tcolorbox}[breakable,enhanced,title=Regiones en el espacio]
\begin{Verbatim}[formatcom=\letraverbatim]
{\[Rho],\[Theta]}={Sqrt[x^2+y^2+z^2],ArcTan[y/Abs[x]]}
\[Phi]=ArcCos[z/\[Rho]];
RegionPlot3D[1<\[Rho]<2&&0<\[Theta]<3Pi/4&&Pi/6<
\[Phi]<3*Pi/4,{x,-2,2},{y,0,2},{z,-2,2},
PlotPoints->90,Mesh->False]
\end{Verbatim}
\end{tcolorbox}

Utilizando la función \verb"ArcTan[x,y]", se puede obtener la región del espacio determinada por 
las expresiones $1\leq \rho  \leq 2,$ $0\leq \theta \leq 3\pi/2,$  $\pi/4\leq \phi \leq 5\pi/6.$


\begin{figure}[H]
\centering
\caption{\small Región determinada por $1\leq\rho\leq 2,$ $0\leq\theta\leq 3\pi/2,$ $\pi/4\leq\phi\leq 5\pi/6.$}
\includegraphics[width=0.4\textwidth]{Img/T_50a.pdf} \hspace{5mm}
\includegraphics[width=0.4\textwidth]{Img/T_50b.pdf}
\fuentepropia
\end{figure}

La figura se consigue a partir de dos grupos 
de desigualdades que se concatenan con el operador lógico \emph{o}, que se denota por \verb"||".
Es decir, la desigualdad $0\leq \theta \leq 3\pi/2,$ se escribe en \verb"Mathamatica" como 
$0\leq \theta \leq \pi/2 \, || \,-\pi \leq \theta \leq -\pi/2.$  
 
\begin{tcolorbox}[breakable,enhanced,title=Instrucciones en Mathematica]
\begin{Verbatim}[formatcom=\letraverbatim]
{\[Theta],\[Rho]}={ArcTan[x,y],Sqrt[x^2+y^2+z^2]}
\[Phi]=ArcCos[z/\[Rho]]; 
RegionPlot3D[1<\[Rho]<2&&0<\[Theta]<Pi&&Pi/4<
\[Phi]<5*Pi/6||1<\[Rho]<2&&-Pi<\[Theta]<-Pi/2
&&Pi/4<\[Phi]<5*Pi/6,{x,-2,2},{y,-2,2},{z,-2,2},
PlotPoints->90,Mesh->False] 
\end{Verbatim}
\end{tcolorbox}
 
La función \verb"ToSphericalCoordinates[{x,y,z}]" permite transformar coordenadas rectangulares en esféricas, mientras que la transformación inversa se logra con \Verb!FromSphericalCoordinates[{\[Rho],\[Theta],\[Phi]}]!.
Ejemplo de lo anterior son la siguientes sentencias:
\begin{tcolorbox}[breakable,enhanced]
\begin{Verbatim}[formatcom=\letraverbatim]
ToSphericalCoordinates[{1,-1,1}]
FromSphericalCoordinates[{Sqrt[3],ArcTan[1],Pi/4}]
\end{Verbatim}
\end{tcolorbox}
\index{ToSphericalCoordinates} \index{FromSphericalCoordinates} 

\subsection*{Gráficos con restricciones}
Agregando \verb"RegionFunction->Function[{x,y,z},DESIGULADADES]" \\
dentro de las opciones
de un gráfico se pueden agregar restricciones a este. \index{RegionFunction}
Sobre la esfera $x^2+y^2+z^2=9$ se aplica la restricción $-2yz<x^2+z^2<9$
y el paraboloide $z=5-x^2-y^2$ con la condición $2x<x^2+y^2<9$.


\begin{figure}[H]
\centering
\caption{\small Superficies con restricciones en el espacio.}
\includegraphics[width=0.4\textwidth]{Img/T_47a.pdf} \hspace{5mm}
\includegraphics[width=0.4\textwidth]{Img/T_48a.pdf}
\fuentepropia
\end{figure}

Estas figuras se obtienen con las siguientes instrucciones:
%

\begin{tcolorbox}[breakable,enhanced,title=Uso del comando {\tt RegionFunction} ]
\begin{Verbatim}[formatcom=\letraverbatim]
ContourPlot3D[x^2+y^2+z^2==9,{x,-3,3},{y,-3,3},
{z,-3,3},RegionFunction->Function[{x,y,z},
-2*y*z<x^2+z^2<9],Mesh->False]
Plot3D[5-x^2-y^2,{x,-3,3},{y,-3,3},Mesh->False,
RegionFunction->Function[{x,y,z},2*x<x^2+y^2<9]]
\end{Verbatim}
\end{tcolorbox}

\subsection*{Gráficos discretos en el plano}
Con la instrucción \Verb!DiscretePlot[f(x),{x,a,b,r},ExtentSize->Right]! se muestra el gráfico 
discreto de $f(x)$ en $[a,b]$, con subintervalos de longitud $r.$ La disposición (o altura) de cada rectángulo
se controla con \verb"ExtentSize" cuyas opciones son \Verb!Left/Right/None/Full/All/Scaled["$\Delta$\verb"]!; 
también se puede incluir la opción \Verb!ExtentSize->{a,b}!, en donde \verb"a" y \verb"b" son la distancia del punto 
y el ancho del rectángulo.


En la figura 25 algunos puntos sobre la función $y=4-x+3x^2-x^3$, $x\in [-1,3]$ e 
incrementos de $0.2$ unidades. 


\begin{figure}[H]
\centering
\caption{\small Función $f(x)=4-x+3x^2-x^3$.}
\includegraphics[width=0.4\textwidth]{Img/T_43b.pdf} \hspace{5mm}
\includegraphics[width=0.4\textwidth]{Img/T_43a.pdf}
\fuentepropia
\end{figure}

Además de estas, las siguientes instrucciones muestran otras posibilidades.

\begin{tcolorbox}[breakable,enhanced,title=Instrucciones en Mathematica]
\begin{Verbatim}[formatcom=\letraverbatim]
DiscretePlot[4-x+3*x^2-x^3,{x,-1,3,0.2}]
DiscretePlot[4-x+3*x^2-x^3,{x,-1,3,0.2},
ExtentSize->All]
Show[DiscretePlot[4-x+3*x^2-x^3,{x,-1,3,0.2},
ExtentSize->All],Plot[4-x+3*x^2-x^3,{x,-1,3}]]
Show[DiscretePlot[4-x+3*x^2-x^3,{x,-1,3,0.24},
ExtentSize->{0.1,0.1}],Plot[4-x+3*x^2-x^3,
{x,-1,3},PlotStyle->Black]]
\end{Verbatim}
\end{tcolorbox}

\subsection*{Gráficos discretos en el espacio}
La instrucción \verb"ListPoint" permite graficar una lista de puntos en el plano, mientras que \verb"ListPointPlot3D" permite hacerlo en el espacio.
 
\index{ListPoint} \index{ListPointPlot3D}
Las siguientes figuras corresponden a la función $f(x,y)= 2+xy-x^2-y^3.$ Para esto, se genera una tabla con puntos para 
$x \in[-2,2],$ $y \in[-2,2]$ con incrementos en $x$ de $1/5$ e incrementos en $y$ de $1/4.$ 

\begin{figure}[H]
	\begin{center}
		\caption{Función $f(x,y)=2+xy-x^2-y^3$.}
		\includegraphics[height=4.5cm]{Img/T_44a.pdf}
		\includegraphics[height=4.5cm]{Img/T_44b.pdf}
		\includegraphics[height=4.5cm]{Img/T_44c.pdf}
	\end{center}
	\fuentepropia
\end{figure}


Las instrucciones necesarias para obtener la figura 26 son las siguientes: en ellas 
se usó el comando \verb"InterpolationOrder" en dos posibilidades.
\begin{tcolorbox}[breakable,enhanced,title=Instrucciones en Mathematica]
\begin{Verbatim}[formatcom=\letraverbatim]
ListPointPlot3D[Table[2+i*j-i^2-j^3,{i,-2,2,1/5},
{j,-2,2,1/4}]]
ListPlot3D[Table[2+i*j-i^2-j^3,{i,-2,2,1/5},
{j,-2,2,1/4}],
Mesh->None,InterpolationOrder->0]
ListPlot3D[Table[2+i*j-i^2-j^3,{i,-2,2,1/5},
{j,-2,2,1/4}],Mesh->None,InterpolationOrder->2]
\end{Verbatim}
\end{tcolorbox}

A partir de la función $z=f(x,y)$ con $x \in [a,b]$ e incrementos de $r_1$ unidades, $y \in [c,d]$ 
e incrementos de $r_2$ unidades, se puede trazar la figura de la función discreta 
obtenida mediante el uso de la instrucción 

\begin{tcolorbox}[breakable,enhanced,title=Instrucciones en Mathematica]
\begin{Verbatim}[formatcom=\letraverbatim]
DiscretePlot3D[f[x,y],{x,a,b,r1},{y,c,d,r2}]
\end{Verbatim}
\end{tcolorbox}
\index{DiscretePlot3D}

Considerando la función $f(x,y)=3+xy-x^2-y^3$ con $x \in [-1,1],$ $y \in [-1,1],$ e incrementos de $0.1$ y
$0.2$ respectivamente, se obtiene los siguientes gráficos.

 % Asegura alineación con el margen izquierdo


\begin{figure}[H]
\centering
\caption{\small Función $f(x,y)=3+xy-x^2-y^3$.}
\label{fig:funcion_T45}
\includegraphics[width=0.4\textwidth]{Img/T_45a.pdf} \hspace{5mm}
\includegraphics[width=0.4\textwidth]{Img/T_45b.pdf}
\end{figure}

Los gráficos anteriores se obtienen al ejecutar las siguientes sentencias.


\begin{tcolorbox}[breakable,enhanced,title=Instrucciones en Mathematica]
\begin{Verbatim}[formatcom=\letraverbatim]
DiscretePlot3D[3+x*y-x^2-y^3,{x,-1,1,0.2},{y,-1,1,0.1},
ExtentSize->Full]
DiscretePlot3D[3+x*y-x^2-y^3,{x,-1,1,0.2},{y,-1,1,0.2},
ExtentSize->{0.15,0.18},FillingStyle->Opacity[0.4]]
\end{Verbatim}
\end{tcolorbox}

\index{FillingStyle}
Los rectángulos se han separan un poco y con la instrucción \verb"FillingStyle" se hacen un poco traslúcidos.

Se presenta enseguida la grafica discreta de la función $f(x,y)=3+xy-x^2-y^3$ 
con la restricción $\frac{1}{4} \leq x^2+y^2 \leq 2$;
e incrementos en $x$ e $y$ de $0.1$ unidades en ambos casos.


\begin{figure}[H]
\centering
\caption{\small Función $f(x,y)=3+xy-x^2-y^3$.}
\includegraphics[width=0.4\textwidth]{Img/T_46a.pdf} \hspace{5mm}
\includegraphics[width=0.4\textwidth]{Img/T_46b.pdf}
\fuentepropia
\end{figure}

Este gráfico se consigue al ejecutar las siguientes instrucciones. La condición se obtiene al usar la opción \verb"RegionFunction". \index{RegionFunction}

\begin{tcolorbox}[breakable,enhanced,title=Instrucciones en Mathematica]
\begin{Verbatim}[formatcom=\letraverbatim]
DiscretePlot3D[3+x*y-x^2-y^3,{x,-2,2,0.1},{y,-2,2,0.1},
ExtentSize->{0.1,0.1},FillingStyle->Opacity[0.6],
RegionFunction->Function[{x,y,z},1/4<=x^2+y^2<=2]]
\end{Verbatim}
\end{tcolorbox}

\subsection*{Superficie de revolución}
El gráfico de la región que se obtiene de girar la función 
$y=f(x)$ alrededor del eje $y$, con $x\in[a,b]$, se obtiene con ayuda 
de la siguiente instrucción:
\begin{tcolorbox}[breakable,enhanced,title=Sólidos de revolución] 
\verb"RevolutionPlot3D[f(x),{x,a,b}]"
\end{tcolorbox}  
\index{RevolutionPlot3D}
Si la curva está parametrizada en la forma $\langle x(t),y(t),z(t) \rangle$ y esta gira alrededor del vector $\langle p,q,r \rangle,$ con  $t \in[a,b]$, la sintaxis que
se usa es la siguiente:
\begin{tcolorbox}[breakable,enhanced,title=Sólido de revolución con una curva parametrizada]
\verb"RevolutionPlot3D[{x(t),y(t),z(t)},{t,a,b},"\\
\verb"RevolutionAxis->{p,q,r}]"
\end{tcolorbox}
Por ejemplo, la función 
${\bf r}(t)= \langle 1 - \cos(t), \sin(t), t - t^2\rangle $, con $t\in [0, 1]$, gira alrededor del vector 
$\langle  1,1,0\rangle$.


\begin{figure}[H]
\centering
\caption{\small Superficie generada por ${\bf r}(t)= \langle 1 - \cos(t), \sin(t), t - t^2\rangle.$}
\includegraphics[width=0.4\textwidth]{Img/math_06ad.pdf} \hspace{5mm}
\includegraphics[width=0.4\textwidth]{Img/math_06ae.pdf}
\fuentepropia
\end{figure}

Esta región se obtiene al ejecutar la siguiente instrucción:
\begin{tcolorbox}[breakable,enhanced] 
\begin{Verbatim}[formatcom=\letraverbatim]
RevolutionPlot3D[{1-Cos[t],Sin[t],t-t^2},
{t,0,1},Ticks->{{0,0.5},{0,0.5},{-0.2,0.2}},
RevolutionAxis->{1,1,0}]
\end{Verbatim}
\end{tcolorbox}

\section{Regresión e interpolación}
A partir de un conjunto de puntos es posible obtener una función
que los contenga o los aproxime, presentaremos algunas funciones que permiten hacerlo.
 
\textls[-10]{Para obtener una combinación lineal $f(x),$ de las funciones 
$f_{1},f_{2},\cdots ,f_{k},$ que mejor ajusta a los puntos $\left( x_{1},y_{1}\right) ,\left(
x_{2},y_{2}\right) ,\cdots ,\left( x_{n},y_{n}\right)$ se puede usar la función \verb"Fit".  \index{Fit}
La función $f(x)$ minimiza la siguiente suma de los cuadrados}  
\begin{multline*}
\left( a_{1}f_{1}\left( x_{1}\right) +a_{2}f_{2}\left( x_{1}\right) 
 +\cdots +a_{n}f_{n}\left( x_{1}\right) -y_{1}\right) ^{2}+ \\ 
 \cdots +\left(a_{1}f_{1}\left( x_{1}\right) +a_{2}f_{2}\left( x_{1}\right) +
\cdots +a_{n}f_{n}\left( x_{1}\right) -y_{n}\right) ^{2}
\end{multline*}
En este caso se puede utilizar la siguiente instrucción genérica:
\begin{tcolorbox}[breakable,enhanced,title=Ajuste de puntos]
\begin{Verbatim}[formatcom=\letraverbatim]
Fit[{{x1,y1},{x2,y2},...,{xn,yn}},{f1,f2,...,fk},x].
\end{Verbatim}
\end{tcolorbox}
Si el usuario establece la forma de la función que desea, los coeficientes de esta combinación 
se consiguen con la función \verb"FindFit".  \index{FindFit}
 
La combinación lineal de las funciones $1,$ $x^2,$ $\sin x,$
que mejor ajusta los puntos  $(-1, 1),$ $(2,-1)$, $(1,3)$ y $(3,0)$,
se obtiene al ejecutar:
\begin{tcolorbox}[breakable,enhanced]
\begin{Verbatim}[formatcom=\letraverbatim]
Fit[{{-1, 1},{2,-1},{1,3},{3,0}},{1,x^2,Sin[x]},x]
\end{Verbatim}
\end{tcolorbox}
Si se desea obtener una expresión de la forma $f(x)=ax+bx^2+c \sin x$, 
los coeficientes se determinan con la función \verb"FindFit," esto es:
\begin{tcolorbox}[breakable,enhanced]
\begin{Verbatim}[formatcom=\letraverbatim]
FindFit[{{1,1},{2,2},{-1,2},{-2,2}},{a*x+b*x^2+c*Sin[x]},
{a,b,c},x].
\end{Verbatim}
\end{tcolorbox}
En las siguientes tres sentencias se define un conjunto de puntos, se halla la
combinación lineal de las funciones $1,x^{2},x^{3},x^{5}$ que mejor
ajustan estos puntos y se presenta su gráfico.
\begin{tcolorbox}[breakable,enhanced]
\begin{Verbatim}[formatcom=\letraverbatim]
datos={{1,1},{2,2},{-1,2},{-2,2}};
y1=Fit[datos,{1,x^2,x^3,x^5},x]
Show[Graphics[{Red,Point[datos]}],Plot[y1,{x,-2,2}]]
\end{Verbatim}
\end{tcolorbox}
Los siguientes tres ejemplos ilustran el uso de esta función dos y tres
variables independientes:

\begin{tcolorbox}[breakable,enhanced,title=Instrucciones en Mathematica]
\begin{Verbatim}[formatcom=\letraverbatim]
Fit[{{-1,1,2},{2,2,0},{-1,3,3}},{x^2*y,x*y^2},{x,y}]
Fit[{{0,Pi/2,2},{Pi/3,Pi/2,-1},{2*Pi/3,Pi,-2},
{-Pi,2*Pi,1}},{Sin[x],Cos[y],Cos[x*y]},{x,y}]
Fit[{{1,1,2,2},{1,2,2,0},{-1,-2,3,4},{-3,1,2,-1},
{-2,4,0,1}},{x^2*y*z,x*y^2,y*z^2,Sin[x+y],Cos[y+z]},
{x,y,z}]
\end{Verbatim}
\end{tcolorbox}

Para los puntos $(x_1,y_1),$ $(x_2,y_2),\cdots$ $(x_n,y_n),$ existe un polinomio $p(x)$ 
de grado $n-1$, tal que $p(x_i)=y_i$ para $i=1,2,\cdots,n$. 
Este polinomio se denomina \emph{polinomio de interpolación} y se determina 
de la siguiente manera:
\begin{tcolorbox}[breakable,enhanced,title=Polinomio de interpolación]
\begin{Verbatim}
InterpolatingPolynomial[{{x1,y1},{x2,y2},...
,{xn,yn}},x]
\end{Verbatim}
\end{tcolorbox}
\index{InterpolatingPolynomial}

Con las siguientes instrucciones se define un conjunto de puntos, se halla
el polinomio de interpolación y se presenta su gráfico
\begin{tcolorbox}[breakable,enhanced]
\begin{Verbatim}[formatcom=\letraverbatim]
puntos={{-1,1},{0,-1},{2,1},{4,-2}};
y=InterpolatingPolynomial[puntos,x];
Show[Graphics[{Red,Point[puntos]}],Plot[y,{x,-1,4}]]
\end{Verbatim}
\end{tcolorbox}

\section{Integración múltiple}
Para evaluar integrales dobles o triples en \verb"Wolfram Mathematica", 
se puede proceder como sigue.\\[0.2cm]
Para una integral doble de la forma 
$\ds \int_a^b \int_c^d f(x,y)dydx,$ la instrucción es 
\begin{tcolorbox}[breakable,enhanced,title=Integrales dobles]
\begin{Verbatim}[formatcom=\letraverbatim]
Integrate[f[x,y],{x,a,b},{y,c,d}]
\end{Verbatim}
\end{tcolorbox}
De manera alternativa, se puede utilizar la siguiente notación:\\
\verb"Integrate[Integrate[f[x,y,z],{x,a,b}],{y,c,d}]"; nótese en este caso el orden en la escritura. 
El comando \verb"NIntegrate" pemite evaluar numéricamente una integral. \index{NIntegrate}
Si se desea ver la integral sin evaluar se puede 
agregar la sentencia \verb"//HoldForm" y si se desea evaluar numéricamente una expresión se debe agregar
\verb"//N" al final de la instrucción. Presentamos algunas integrales y las instrucciones requeridas para su evaluación:
%\clearpage
\begin{enumerate}[leftmargin=0cm,itemindent=.75cm,labelwidth=\itemindent,labelsep=0cm,align=left,label={{\bf(\arabic*)}}]
\item %La integral 
$$\ds\int_0^2 \int_{-x}^{x^2} xy^2 dxdy,$$ %se evalúa con la siguiente orden
\begin{tcolorbox}[breakable,enhanced] 
\begin{Verbatim}
Integrate[x*y^2,{x,0,2},{y,-x,x^2}]
\end{Verbatim}
\end{tcolorbox}
\item La siguiente integral se evalúa numéricamente 
$$\ds \int_0^2 \int_{\sqrt{y}}^{2-y} 30x^2 dxdy,$$ %se puede usar la siguiente instrucción:
\begin{tcolorbox}[breakable,enhanced] 
\begin{Verbatim}[formatcom=\letraverbatim]
NIntegrate[30x^2,{y,0,1},{x,Sqrt[y],2-y}]
\end{Verbatim}
\end{tcolorbox}
\item %La integral impropia 
$$\int_{-\infty}^{\infty}\int_{0}^{\infty }e^{-x^{2}-y^{2}}dydx= \frac{\pi}{2} $$ % se evalúa como sigue
\begin{tcolorbox}[breakable,enhanced] 
\begin{Verbatim}[formatcom=\letraverbatim]
Integrate[Exp[-x^2-y^2],{x,-Infinity,Infinity},
{y,0,Infinity}]
\end{Verbatim}
\end{tcolorbox}
\end{enumerate}


La integral triple $$\int_a^b \int_c^d \int_e^s f(x,y,z)\, dzdydx,$$
se puede evaluar mediante esta instrucción:
\begin{tcolorbox}[breakable,enhanced,title=Integrales triples]
\begin{Verbatim}[formatcom=\letraverbatim]
Integrate[f[x,y,z],{z,r,s},{y,c,d},{x,a,b}]
\end{Verbatim}
\end{tcolorbox}
Ahora ilustraremos el uso de esta función con unas integrales triples:
\begin{enumerate}[leftmargin=0cm,itemindent=.75cm,labelwidth=\itemindent,labelsep=0cm,align=left,label={{\bf(\arabic*)}}]
\item % La integral, 
$$\ds \int_{-3}^3 \int_0^2 \int_0^1  x^2 y^3 z^4 dzdy dx$$ % se evalúa como sigue 
\begin{tcolorbox}[breakable,enhanced] 
\begin{Verbatim}[formatcom=\letraverbatim]
Integrate[x^2y^3z^4,{x,-3,3},{y,0,2},{z,0,1}]
\end{Verbatim}
\end{tcolorbox}
\item % La evaluación de la integral 
$$\ds \int_{0}^{1}\int_{-\sqrt[2]{x}}^{\sqrt[2]{x}}\int_{0}^{x}x^{2}dzdydx$$ % se hace con la orden
\begin{tcolorbox}[breakable,enhanced] 
\begin{Verbatim}[formatcom=\letraverbatim]
Integrate[x^2,{x,0,1},{y,-Sqrt[x],Sqrt[x]},{z,0,x}]
\end{Verbatim}
\end{tcolorbox}

\item Agregando \verb"//HoldForm" al final de la siguiente sentencia, se puede ver la forma sin evaluar de la integral:
$$\int_0^\pi\int_0^{\pi/6} \int_0^2 \rho^3 \cos(\theta)\sin^2(\phi)\, d\rho d \theta d\phi=4  $$
\begin{tcolorbox}[breakable,enhanced]
\begin{Verbatim}[formatcom=\letraverbatim] 
Integrate[\[Rho]^3*Cos[\[Theta]]*Sin[\[Phi]],
{\[Rho],0,2},{\[Theta],0,Pi/6},{\[Phi],0,Pi}]
\end{Verbatim}
\end{tcolorbox}

%\item La integral doble impropia, $$\int_{-\infty }^{\infty }\int_{-\infty }^{\infty }x^{2}e^{-x^{2}-y^{2}}dydx= \frac{\pi }{2} $$ se evalúa mediante la ejecución de la siguiente instrucción 
%\begin{tcolorbox}[breakable,enhanced]  \begin{Verbatim}[fontsize=\letraverbatim]
%Integrate[x^2*Exp[-x^2-y^2],{x,-Infinity,Infinity}, {y,-Infinity,Infinity}]\end{Verbatim}\end{tcolorbox}
\item % La integral triple impropia
$$\frac{1}{\sqrt{\pi }}\int_{-\infty}^{\infty}\int_{-\infty}^{\infty}\int_{-\infty}^{\infty}e^{-x^{2}-y^{2}-z^{2}}dzdydx= \pi $$
% se evalúa mediante la ejecución de la siguiente instrucción 
\begin{tcolorbox}[breakable,enhanced] 
\begin{Verbatim}[formatcom=\letraverbatim]
1/Sqrt[Pi]*Integrate[Exp[-x^2-y^2-z^2],
{x,-Infinity,Infinity},{y,-Infinity,Infinity},
{z,-Infinity,Infinity}]
\end{Verbatim}
\end{tcolorbox}
\end{enumerate}

Muchos de los cálculos y manipulaciones algebraicas de estas notas se hacen en \verb"Mathematica" 
y se exportan a un documento \LaTeX.
El código de una \verb"EXPRESION" se consigue mediante la instrucción
\verb"TeXForm[EXPRESION]". \index{TeXForm}
\section{Ecuaciones diferenciales con {\tt DSolve}}
Se denomina \emph{ecuación diferencial} a toda ecuación que involucre
una función desconocida y algunas de sus derivadas. La ecuación
diferencial se llama ecuación diferencial ordinaria, si contiene
derivadas ordinarias y se llama ecuación diferencial parcial, si en ella
aparecen algunas derivadas parciales de una función que depende de más de una variable. 
El \emph{orden} de la ecuación diferencial corresponde al orden de la mayor derivada involucrada.
Una \emph{solución} de la ecuación diferencial es una función definida en un conjunto abierto $I$, tal que para todo $x\in I$ sus derivadas existen y al reemplazarse en la ecuación diferencial la expresión es una identidad.
Una introducción a las ecuaciones diferenciales puede consultarse en \textcite{apostol_calculus_1999}.
\index{Ecuaciones diferenciales} \index{DSolve}
 El comando \verb"DSolve" resuelve una o varias ecuaciones diferenciales ordinarias o parciales.
La ecuación diferencial $y'(x)=f(x)$ se resuelve en \verb"Mathematica" mediante la ejecución de la siguiente instrucción
\verb"DSolve[y'[x]==f[x],y[x],x]". \\ Según el caso, se debe usar agregar como una lista de ecuaciones o una ecuación con sus respectivas condiciones.
\begin{tcolorbox}[breakable, enhanced, title={Uso del comando \texttt{DSolve}}, colback=gray!5, colframe=gray!80]
El comando \texttt{DSolve} resuelve una o varias ecuaciones diferenciales ordinarias o parciales.
La ecuación diferencial $y'(x) = f(x)$ se resuelve en \texttt{Mathematica} mediante la instrucción:
\texttt{DSolve[y'[x] == f[x], y[x], x]}.
Según el caso, se debe usar una lista de ecuaciones o una ecuación con sus respectivas condiciones.
\end{tcolorbox}

El comando \verb"NDSolve" permite resolver numéricamente una ecuación diferencial,
una exposición más detallada de esta función puede conseguirse \textcite{sofroniou_advanced_2008}. Ilustraremos lo anterior con unos ejemplos.
\begin{itemize}[leftmargin=0cm,itemindent=.5cm,labelwidth=\itemindent,labelsep=0cm,align=left]
\item La ecuación diferencial ordinaria de segundo orden
$$ y^{\prime \prime }\left( x\right) + y^{\prime }\left(x\right) =-e^{-x},$$ con las condiciones
$y\left( 0\right) =1,$ $y^{\prime }\left( 0\right) =1,$ tiene por solución $y\left( x\right) =xe^{-x}+1.$
\begin{tcolorbox}[breakable,enhanced]
\begin{Verbatim}[formatcom=\letraverbatim] 
DSolve[{y''[t]+y'[t]==-Exp[-t],y[0]==1,y'[0]==1},
y[t],t]
\end{Verbatim}
\end{tcolorbox}

\item La ecuación diferencial parcial de primer orden
$$\frac{\partial u}{\partial x}+\frac{\partial u}{\partial y}=e^{-y}\left(1-x\right) ,$$ con $u\left( 0,y\right) =0,$ $u\left( x,0\right) =x,$ posee a $u\left( x,y\right) =xe^{-y}$ como solución.

\begin{tcolorbox}[breakable,enhanced]
\begin{Verbatim}[formatcom=\letraverbatim] 
DSolve[{D[u[x,y],x]+D[u[x,y],y]==Exp[-y]*(1-x),
u[0,y]==0,u[x,0]==x},u[x,y],{x,y}]
\end{Verbatim}
\end{tcolorbox}

 
\item El sistema de ecuaciones diferenciales ordinarias de primer orden
$$\left. \begin{array}{l}
y'(t)-x'(t)+y(t)=-e^t(t+1)\\
x'(t)+y'(t)-x(t)=e^t-e^{-t}
\end{array}\right \}$$
con las condiciones iniciales $x(0)=0$, $y(0)=1$ posee a $x(t)= te^t$ y $y(t)=e^{-t}$ como solución.
 
\begin{tcolorbox}[breakable,enhanced,title=Instrucciones en Mathematica]
\begin{Verbatim}[formatcom=\letraverbatim]
DSolve[{y'[t]-x'[t]+y[t]==-Exp[t]*(t+1),x'[t]
+y'[t]-x[t]==Exp[t]-Exp[-t],x[0]==0,y[0]==1},
{x[t],y[t]},t]
\end{Verbatim}
\end{tcolorbox}

  
\end{itemize}

\section{Campos vectoriales}
Presentamos algunas funciones del programa \verb"Mathematica" útiles en el tratamiento de campos vectoriales.

\subsection*{Divergencia y rotacional} \index{Curl} \index{Div}
Con las funciones \verb"Curl" y \verb"Div" se obtiene el rotacional y la divergencia de un campo vectorial 
${\bf F}.$ Se debe establecer un orden para las variables.
\index{Divergencia} \index{Rotacional} 

En el siguiente ejemplo se calcula el rotacional y la divergencia de ${\bf F}$, \mbox{después} de definir 
el campo vectorial ${\bf F}(x,y,z)= \la x^2+y^2,y^3+z^2,x e^z \ra.$
\begin{tcolorbox}[breakable,enhanced]
\begin{Verbatim}[formatcom=\letraverbatim]
F[x_,y_,z_]:={x^2+y^2,y^3+z^2,x*Exp[z]}
Curl[F[x,y,z],{x,y,z}]
Div[F[x,y,z],{x,y,z}]
\end{Verbatim}
\end{tcolorbox}

\subsection*{Gráficos de campos vectoriales} \index{VectorPlot}
Con la instrucción \verb"VectorPlot[F(x,y),{x,a,b},{y,c,d}]" se obtiene un gráfico del campo vectorial
${\bf F}(x,y)$, para $x \in [a,b]$, $y\in [c,d]$. 
La sintaxis que se debe usar para obtener una representación del campo vectorial ${\bf F}(x,y,z)$ es
\verb"VectorPlot3D[F[x,y,z],{x,a,b},{y,c,d},{z,p,q}]". \index{VectorPlot3D}
El campo vectorial ${\bf F}(x,y)= \la \frac{-x-y}{\sqrt{x^2+y^2}}, \frac{x-y}{\sqrt{x^2+y^2}} \ra$, 
se representa para $x \in [-1,1],$ $y \in [-1,1].$ También se presentan algunas \emph{líneas de flujo}\footnote{Para un campo vectorial ${\bf F}$, una línea de flujo para ${\bf F}$ es una curva ${\bf r}(t)$ tal que
${\bf r}^\prime(t) = {\bf F}({\bf r}(t)).$ Esto es,
${\bf F}$ corresponde al campo de velocidades de la curva ${\bf r}(t)$}, esto se logra
utilizando la función \verb"StreamPlot".  \index{StreamPlot} \index{StreamStyle}\index{StreamScale}

\begin{figure}[H]
	\begin{center}
		\caption{Campo vectorial $\la \frac{-x+y}{\sqrt{x^2+y^2}}, \frac{x-y}{\sqrt{x^2+y^2}} \ra$.}
		\includegraphics[height=4.5cm]{Img/T_55a.pdf} \hspace{0.5cm}
		\includegraphics[height=4.5cm]{Img/T_55b.pdf}
		\parbox{0.95\linewidth}{\centering\fontsize{10pt}{10pt}\selectfont\textbf{Fuente:} elaboración propia.}
	\end{center}
\end{figure}

Las instrucciones requeridas son, respectivamente las siguientes.

\begin{tcolorbox}[breakable,enhanced,title=Instrucciones en Mathematica]
\begin{Verbatim}[formatcom=\letraverbatim] 
VectorPlot[{-(x+y)/Sqrt[x^2+y^2],(x-y)/
Sqrt[x^2+y^2]},
{x,-1,1},{y,-1,1},VectorStyle->Black]
StreamPlot[{-(x+y)/Sqrt[x^2+y^2],(x-y)/
Sqrt[x^2+y^2]},
{x,-1,1},{y,-1,1},
StreamStyle->Black,StreamScale->0.08]
\end{Verbatim}
\end{tcolorbox}

En estas notas nos ocuparemos de obtener la representación gráfica de un campo vectorial ${\bf F}(x,y,z)$ sobre una superficie determinada por la expresión $g(x,y,z)=0$, para $x\in [a,b],$  $y\in [c,d]$, $z\in [e,f]$. Esto se 
logra en \verb"Mathematica" utilizando la función \verb"SliceVectorPlot3D". La estructura general es la siguiente:
\begin{tcolorbox}[breakable,enhanced,title=Campos vectoriales sobre superficies] 
\begin{Verbatim}[formatcom=\letraverbatim]
SliceVectorPlot3D[F[x,y,x],{g[x,y,z]=0},{x,a,b},
{y,c,d},{z,e,f},Opciones]
\end{Verbatim}
\end{tcolorbox} \index{SliceVectorPlot3D}



La figura del campo vectorial ${\bf F}(x,y,z)= \la y,-x,x+z \ra$, sobre la superficie del sólido 
acotado por las figuras de las expresiones $z=2(1-x^2-y^2)$ y $z+1=(x^2+y^2)^2$ se obtiene ejecutando 
las siguientes instrucciones. La intersección está determinada por la ecuación $x^2+y^2=1$,
por este motivo se agregó la condición $x^2+y^2\leq 1$.



\begin{tcolorbox}[breakable,enhanced,title=Instrucciones en Mathematica]
\begin{Verbatim}[formatcom=\letraverbatim]
SliceVectorPlot3D[{y,-x,x+z},{(x^2+y^2)^2==z+1,
z==2*(1-x^2-y^2)},{x,-1,1},{y,-1,1},{z,-1,2.2},
RegionFunction->Function[{x,y,z},x^2+y^2<=1],
VectorScale->{0.18,Scaled[0.3],Automatic},
VectorStyle->Black,
VectorPoints->16,PlotStyle->Orange]
\end{Verbatim}
\end{tcolorbox}


\begin{figure}[H]
\centering
\caption{\small Campo vectorial ${\bf F}(x,y,z)= \la y,-x,x+z \ra$.}
\includegraphics[width=0.4\textwidth]{Img/T_28a.pdf} \hspace{5mm}
\includegraphics[width=0.4\textwidth]{Img/T_28b.pdf}
\end{figure}

La figura del campo vectorial ${\bf F}(x,y,z)= \la -y,x,z/2 \ra$ sobre la superficie del paraboloide $4z=4-x^2-y^2$ con la condición $x^2+y^2\leq 4$. 


\begin{figure}[H]
\centering
\caption{\small Campo vectorial ${\bf F}(x,y,z)= \la -y,x,z/2 \ra$.}
\includegraphics[width=0.4\textwidth]{Img/T_27a.pdf} \hspace{5mm}
\includegraphics[width=0.4\textwidth]{Img/T_27b.pdf}
\end{figure}

Esta figura se obtiene compilando la siguiente instrucción. En esta se usó el comando \verb"RegionFunction" que restringe el dominio en la figura de la superficie. \index{RegionFunction}



\begin{tcolorbox}[breakable,enhanced,title=Instrucciones en Mathematica]
\begin{Verbatim}[formatcom=\letraverbatim]
SliceVectorPlot3D[{-y,x,z/2},{4*z==4-x^2-y^2},
{x,-2,2},
{y,-2,2},{z,0,1.1},RegionFunction->Function[{x,y,z},
x^2+y^2<=4],
VectorScale->{0.15,Scaled[0.25],Automatic},
VectorStyle->Black,VectorPoints->16,
PlotStyle->Orange,
Ticks->{{-2,-1,0,1,2},{-2,-1,0,1,2},{0,1}}]
\end{Verbatim}
\end{tcolorbox}

También se puede obtener la representación gráfica de un campo vectorial con
la opción \verb"ListSliceVectorPlot3D." Si se considera el campo vectorial ${\bf F}(x,y,z)$
para $x\in[a,b],$ $y\in[c, d],$ $z\in[e,f]$, se crea una tabla que denominaremos datos,
posteriormente se dibuja el campo vectorial sobre la superficie determinada por
la expresión $g(x,y,z)=0$. En este caso la sintaxis básica es la siguiente:

\enlargethispage{\baselineskip}

\begin{tcolorbox}[breakable,enhanced,title=Instrucciones en Mathematica]
\begin{Verbatim}[formatcom=\letraverbatim] 
datos=Table[F[x,y,z],{x,a,b},{y,c,d},{z,e,f}];
ListSliceVectorPlot3D[datos,
ImplicitRegion[g[x,y,z]==0,{x,y,z}],
DataRange->{{a,b},{c,d},{e,f}}]
\end{Verbatim}
\end{tcolorbox}

El campo vectorial ${\bf F}(x,y,z) =\la x+z,y,2x-z \ra$ se representa 
sobre la superficie $z= (y^2-x)e^{1-x^2-y^2}$. 


\begin{figure}[H]
\centering
\caption{\small Campo vectorial ${\bf F}(x,y,z) =\la x+z,y,2x-z \ra$.}
\includegraphics[width=0.4\textwidth]{Img/T_33a.pdf} \hspace{5mm}
\includegraphics[width=0.4\textwidth]{Img/T_33b.pdf}
\end{figure}

Las instrucciones requeridas para obtener esta figura son las siguientes:
%ContourPlot3D[z==(y^2-x)*Exp[1-x^2-y^2],{x,-2,2},{y,-2,2},
%{z,-2,2},MeshStyle->Gray,ContourStyle->Orange]

\begin{tcolorbox}[breakable,enhanced,title=Instrucciones en Mathematica]
\begin{Verbatim}[formatcom=\letraverbatim] 
datos=Table[{x+z,y,2*x-z},{x,-2,2},{y,-2,2},
{z,-1,1}];
ListSliceVectorPlot3D[datos,
ImplicitRegion[z==(y^2-x)*Exp[1-x^2-y^2],{x,y,z}],
VectorPoints->12,DataRange->{{-2,2},{-2,2},{-2,2}},
VectorStyle->{Thin,Black,Thick}]
\end{Verbatim}
\end{tcolorbox}

La representación del campo vectorial ${\bf F}(x,y,z) =\la x-z,y,z+2 \ra$ sobre la superficie
del paraboloide hiperbólico $4z=x^2-y^2$ con la condición $x^2+y^2 \leq 4$, se obtiene con las 
siguientes instrucciones:

\begin{tcolorbox}[breakable,enhanced,title=Instrucciones en Mathematica]
\begin{Verbatim}[formatcom=\letraverbatim] 
datos=Table[{x-z,y,z+2},{x,-2,2},{y,-2,2},{z,-1,1}];
ListSliceVectorPlot3D[datos,
ImplicitRegion[4*z==x^2-y^2&&x^2+y^2<=4,{x,y,z}],
VectorPoints->12,DataRange->{{-2,2},{-2,2},
{-5/4,5/4}},VectorStyle->{Thin,Black,Thick}]
\end{Verbatim}
\end{tcolorbox}


\begin{figure}[H]
\centering
\caption{\small Campo vectorial ${\bf F}(x,y,z) =\la x-z,y,z+2 \ra$.}
\includegraphics[width=0.4\textwidth]{Img/T_54a.pdf} \hspace{5mm}
\includegraphics[width=0.4\textwidth]{Img/T_54b.pdf}
\end{figure}

\subsection*{Campos conservativos y función potencial}
Recordemos que si ${\bf F} = P(x, y,z)\vi +Q(x,y,z)\vj +P(x,y,z)\vk$ 
es un campo vectorial conservativo existe un campo escalar $\phi $ tal que $\nabla \phi ={\bf F}.$ 

%\clearpage
\begin{ejemplo}
El campo vectorial 
{\small ${\bf F}(x,y,z) = \left\langle y^{3}e^{-z},3xy^{2}e^{-z},3z^{2}-xy^{3}e^{-z}\right\rangle,$}
es conservativo. Las funciones $P$, $Q$ y $R$ son respectivamente las funciones componentes de ${\bf F}$. 
Para hallar la función potencial \index{Función!potencial} proceder como sigue, 
se integra $P$ con respecto a $x$, esto es 
$$\phi (x,y,z) = \int \left( y^{3}e^{-z}\right)
dx= xy^{3}e^{-z}+g\left( y,z\right) $$
Ahora se compara $\frac{\partial \phi }{\partial y}$ con $Q\left(x,y,z\right) ,$ esto es
$\frac{\partial \phi }{\partial y}= 3xy^{2}e^{-z}+\frac{\partial g}{\partial y}=3xy^{2}e^{-z},$ entonces 
$\frac{\partial g}{\partial y}=0,$ es decir, $g$ no depende de $y,$ por lo tanto, 
$$\phi (x,y,z) = xy^{3}e^{-z}+g\left(z\right) .$$
Luego se compara $\frac{\partial \phi }{\partial z}$ con 
$R(x,y,z),$ esto es $\frac{\partial \phi }{\partial z}= -xy^{3}e^{-z}+g^{\prime}(z) = 3z^{2}-xy^{3}e^{-z}$
entonces $g^{\prime}\left( z\right) =3z^{2},$ de lo cual se sigue que $g\left( z\right) =z^{3}+c.$ Por lo tanto, 
$$\phi (x,y,z) = xy^{3}e^{-z}+z^{3}+c.$$
\end{ejemplo}

En \verb"Wolfram Matematica" este campo escalar puede hallarse mediante la solución de
un sistema de ecuaciones diferenciales parciales: 
$$\left. \begin{array}{c}
\frac{\partial \phi }{\partial x}(x,y,z)=P(x,y,z) \\[0.15cm] 
\frac{\partial \phi }{\partial y}(x,y,z) =Q(x,y,z) \\[0.15cm] 
\frac{\partial \phi }{\partial z}(x,y,z) =R(x,y,z) 
\end{array}%
\right\} $$


Con las siguientes instrucciones además de determinar si el campo vectorial es conservativo, se halla la función potencial.
Se utiliza la función \verb"If", la cual permite elegir entre dos resultado dependiendo si la prueba es verdadera o falsa. \index{If}

\begin{tcolorbox}[breakable,enhanced,title=Instrucciones en Mathematica]
\begin{Verbatim}[formatcom=\letraverbatim]
{P,Q}={y^3*Exp[-z],3x*y^2*Exp[-z]};
R=3*z^2-x*y^3*Exp[-z];
If[D[P,y]==D[Q,x]&&D[P,z]==D[R,x]&&D[Q,z]==D[R,y],
"El Campo vectorial es conservativo",
"El Campo Vectorial No es conservativo"]
DSolve[{D[\[Phi][x,y,z],x]==P,D[\[Phi][x,y,z],y]==Q,
D[\[Phi][x,y,z],z]==R},\[Phi][x,y,z],x,y,z]
\end{Verbatim}
\end{tcolorbox}
\newpage